\section{The Semantics of Generalized Controls}
%
\label{Sect:Semantics}

This section introduces an operation $( \odot )$ called the \emph{control product} which provides a symmetric generalization of controlled operations in unitary quantum circuits.
We begin by defining an assignment $\odot : \CatUnitary \times \CatUnitary \to \CatUnitary$ on unitary matrices.
If $U$ and $V$ are two unitary matrices, then we define $U \odot V = \exp( \Log(U) \otimes \Log(V) / ( i\pi ) )$.
The normalization factor $1 / ( i\pi )$ restricts the eigenvalues of the Hermitian matrix to $( -\pi, \pi ]$, and consequently $( \odot )$ is associative.

\begin{theorem}
    \label{Thm:ControlAssoc}
    The control product is unital and associative with unit $( -1 )$.
    Moreover, given a sequence $U_1, U_2, \ldots, U_n \in \Mor( \CatUnitary )$, $\bigodot_{j=1}^n U_j = \exp\left( i\pi \cdot \bigotimes_{j=1}^{n} \frac{\Log( U_j )}{i\pi} \right)$.
\end{theorem}

Unfortunately, the control product is not a bifunctor since $U \odot ( V \circ W ) \ne ( U \odot V ) \circ ( U \odot W )$ in general.
However, the control product does enjoy many desirable properties.
For example, if the first coordinate is fixed to $Z$, then for each $U \in \CatUnitary( n, n )$, $Z \odot U = \ket{0} \bra{0} \otimes \mathbb{I}_n + \ket{1} \bra{1} \otimes U$.
That is to say, $Z \odot( - )$ is the standard canonical control functor on $\CatUnitary_2$.
It follows that $Z \odot X$ is a \textsf{CNOT} gate and $Z \odot Z$ is a \textsf{CZ} gate.
Since $X \odot ( - )$ is also a control functor, then this formalization suggests that the \textsf{CNOT} gate is also a $Z$-gate controlled in the $X$-basis.
This insight can be used to explain phase kickback~\cite{OssorioCastillo2023}.
Of course, it is natural to ask for which $U \in \Mor( \CatUnitary )$ will $U \odot ( - )$ be a control functor.
This is answered by the following theorem.

\begin{theorem}
    \label{Thm:GenCtrlNF}
    Let $H$ be an $n$-dimensional Hermitian matrix with eigenvalues $\{ \alpha_j \}_{j=1}^n$ in $( -1, 1 ]$ and orthonormal eigenbasis $\{ b_j \}_{j=1}^n$.
    If $U = \exp( iH\pi )$, then $U \odot V = \sum_{j=1}^n \ket{b_j} \bra{b_j} \otimes V^{\alpha_j}$ for each $V \in \Mor( \CatUnitary )$.
    In particular, $Z \odot V = \ket{0} \bra{0} \otimes \mathbb{I}_m + \ket{1} \bra{1} \otimes V$ for $V \in \CatUnitary( m, m )$.
\end{theorem}

\begin{corollary}
    \label{Cor:GenCtrlFunc}
    If $U \in \CatUnitary( d, d )$ and $F( - ) = U \odot ( - )$, then the following conditions hold.
    \begin{itemize}
    \item \textbf{C1}.
          $F( \mathbb{I}_n ) = \mathbb{I}_{dn}$ for each $n \in \mathbb{N}$.
    \item \textbf{C2}.
          $F( V \otimes \mathbb{I}_n ) = F( V ) \otimes \mathbb{I}_n$ for each $n \in \mathbb{N}$ and $V \in \Mor( \CatUnitary )$.
    \item \textbf{C3}.
          $F( F( V ) ) \circ ( \sigma_{d,d} \otimes \mathbb{I}_n ) = ( \sigma_{d,d} \otimes \mathbb{I}_n ) \circ F( F( V ) )$ for each $V \in \CatUnitary( n, n )$ where $\sigma_{d,d}$ is the monoidal symmetry $\mathbb{C}^d \otimes \mathbb{C}^d \cong \mathbb{C}^d \otimes \mathbb{C}^d$.
    \item \textbf{C4}.
          $F( P^\dagger \circ V \circ P ) = ( \mathbb{I}_d \otimes P^\dagger ) \circ F( V ) \circ ( \mathbb{I}_d \otimes P )$ for each $V, P \in \CatUnitary( n, n )$.
    \end{itemize}
    The assignment $F( - )$ is a functor if and only if $U$ is Hermitian.
    Moreover, if $F( - )$ is a functor, then $F( - )$ restricts to a conjugated control functor on $\CatUnitary_d$.
    This control functor is $( u, v )$-pointed if and only if $U\ket{u} = -\ket{u}$ and $U\ket{v} = \ket{v}$.
\end{corollary}

\cref{Thm:GenCtrlNF} establishes a relation between control products and matrix powers.
Moreover, if $\alpha \in (-1, 1]$, then $U^\alpha = U \odot e^{i\alpha\pi}$.
Unfortunately, matrix powers are only well-behaved for small choices of $\alpha$, as illustrated by the following theorem.
In general, the properties of exponentiation will depend on the eigenvalues of $U$.
To this end, we will say that $\lambda_j$ is \emph{$( \alpha, \beta )$-admissible} if either $\alpha \lambda_j \in ( -1, 1  ]$ or $\beta \in \mathbb{Q}$ with reduced denominator dividing $\ceil*{( \alpha \lambda_j - 1) / 2}$.

\begin{theorem}
    \label{Thm:PowAsCtrl}
    For $\alpha \in \mathbb{R}$, $U \odot e^{i\alpha\pi} = U^\alpha$ for all $U \in \Mor( \CatUnitary )$ if and only if $\alpha \in ( -1, 1 ]$.
    If $U \in \Mor( \CatUnitary )$ is Hermitian, then $U \odot e^{i\alpha\pi} = U^\alpha$ for all $\alpha \in \mathbb{R}$.
\end{theorem}

\begin{theorem}
    \label{Thm:AddPow}
    Let $U \in \CatUnitary( n, n )$ with $\{ \lambda_j \}_{j=1}^{n}$ the eigenvalues of $\Log( U ) / (i\pi)$.
    For each $\alpha \in \mathbb{R}$ and $\beta \in \mathbb{R}$, $( U^\alpha )^\beta = U^{\alpha \beta}$ if and only if for each $j \in \{ 1, 2, \ldots, n \}$, $\lambda_j$ is $( \alpha, \beta )$-admissible.
    This means that $\mathbb{Z}$ and $( -1, 1 ]$ are maximal submonoids of $( \mathbb{R}, \cdot )$ for which $( U, \alpha ) \mapsto U^\alpha$ defines a monoid action on $\Mor( \CatUnitary )$.
\end{theorem}

It was shown in~\cite{Schober2024} that generalized controls emerge naturally when studying quantum algorithms using the techniques of~\cite{Gidney2017}.
Moreover, \cref{Thm:GenCtrlNF} and \cref{Thm:PowAsCtrl} show that generalized controls are inherently connected with the matrix powers already used in quantum programming languages such as OpenQASM 3~\cite{Cross2022}.
Ideally, one would hope that the intuition for standard controls and scalar powers would generalize to this setting.
It is clear from \cref{Cor:GenCtrlFunc} and \cref{Thm:AddPow} that neither of these hopes hold true.
For example, integer exponentiation (i.e., repeating and inverting circuits) does not interact well with unit fraction exponentiation (i.e., computing roots of circuits), despite these two operations being fundamentally related.
Given these limitations, one is left to seek out techniques to manipulate control products and matrix powers, even in the cases where the standard intuition fails (e.g., when $\alpha > 1$).
This motivates the study of complete equational theories, which could be used to work soundly with these constructs in the context of circuit optimization, simulation, and equivalence checking.

\begin{example}[Circuit Exponentiation in IBM Heron]
    \label{Ex:RZZ}
    Matrix powers are present in many quantum programming languages, whether explicitly or implicitly, in the form of parametrized rotation gates.
    With the advent of new quantum hardware, multi-qubit rotation gates have been gaining attention.
    For example, the IBM Heron r3 qubit processor released in 2025~\cite{IBMProcessors} offers a multi-qubit $ZZ$-rotation~\cite{IBMRZZ} in addition to the more standard one-qubit $X$- and $Z$-rotation gates, which is intended to reduce execution times and error rates when running quantum simulations and variational algorithms~\cite{IBMFractional}.
    The semantics of IBM Heron r3's parametrized gates are as follows. On the right we include their equivalent semantics in terms of matrix powers.
    \begin{align*}
        R_X(\theta) &:= \exp( -i \theta X / 2 ) = X^{\theta / ( 2\pi )} \cdot (-X)^{-\theta / ( 2\pi )} \\
        R_Z(\theta) &:= \exp( -i \theta Z / 2 ) = Z^{\theta / ( 2\pi )} \cdot (-Z)^{-\theta / ( 2\pi )} \\
        R_{ZZ}(\theta) &:= \exp( -i \theta Z \otimes Z / 2 ) = e^{-i\theta/2} ( Z \otimes Z )^{ \theta / \pi}
    \end{align*}
    A sound and complete equational theory for circuit exponentiation and generalized control would be able to prove the following decomposition of $R_{ZZ}( \theta )$, without the need for new axioms specific to the $ZZ$-rotation.
    \begin{equation*}
        R_{ZZ}( \theta )
        =
        \sem*{\scalebox{0.65}{\circimg{tikz-cache/rzz__decomp_a.pdf}}}
        =
        \sem*{\scalebox{0.65}{\circimg{tikz-cache/rzz__decomp_b.pdf}}}
    \end{equation*}
\end{example}
